\section{Conclusion}

We have rewritten this paper around the result that is actually supported.
Within the five-mode spectral inverse problem for a fluid-filled shell, the
equivalent-sphere reduction is singular in exact arithmetic and appears
numerically as
$\kappa_\mathrm{sphere}\approx\num{1.37e10}$.  Retaining the oblate curvature
field produces a finite near-spherical floor
$\kappa_\mathrm{floor}\approx\num{269}$ because $\sigma_0>0$, and the canonical
oblate geometry improves that further to
$\kappa_\mathrm{oblate}=\num{69.4}$.  A prolate continuation does not provide a
matching benefit, remaining at
$\kappa_\mathrm{prolate}=\numrange{561}{729}$.

The conclusion is therefore specific rather than universal: identifiability in
this spectral inverse problem is geometry-dependent.  Oblate curvature
redistribution lifts the weak inverse direction; prolate distortion does not
deliver the same improvement.  That is a useful result, and unlike the abandoned
universal theorem, it survives contact with the calculations.
